\subsection{Changing the scalar coefficient complex}

Finding~\ref{f:boundedSoergelNoGo} fixes the cohomologically
degree-zero coefficient object. An arbitrary bounded scalar
complex with the same Euler class is a larger premise. The
following necessary conditions and candidate tests address
that premise without extending the preceding theorem beyond
its hypotheses.

Let \(k\) have characteristic zero, let \(n\geq2\), and let
\(A\) be the standard degree-two polynomial realization of
\(W=S_n\times S_n\). Fix a simple reflection \(s\) in both
factors and put \(\tau=(s,s)\). Retain
\(\delta=v+v^{-1}\).
Let \(X\) be a bounded ordinary Soergel complex with
\begin{equation}\label{eq:alteredScalarAmbientClass}
 [X]=P_s=b_s\otimes b_s+
                    (b_s-\delta)\otimes(b_s-\delta).
\end{equation}
Let \(C\) be a bounded complex of finite graded
\(k\)-vector spaces with
\begin{equation}\label{eq:alteredScalarEuler}
 \sum_{j,e}(-1)^j\dim_k H^j(C)_e\,v^e=\delta^2.
\end{equation}
Its coefficient bimodule is \(C\otimes_k A\), with the
ordinary central unit action. Test an actual derived
isomorphism
\begin{equation}\label{eq:alteredScalarQuadratic}
 X\otimes_A^{\mathbf L}X\simeq C\otimes_k X.
\end{equation}
A homotopy quadratic implies this derived quadratic.
Here scalar means a complex over \(k\): a complex of sums
of unit bimodules with nonconstant polynomial differentials
is not automatically of this form.

\begin{finding}[\fid{F59}, altered scalar coefficients]
\label{f:alteredScalarSoergel}
Under \eqref{eq:alteredScalarAmbientClass}--
\eqref{eq:alteredScalarQuadratic}, generic cohomology,
with internal grading forgotten, is supported only on the
identity and \(\tau\) graphs. In each cohomological degree
their multiplicities agree. Precisely, if
\[
 d(u)=\sum_j\dim_k H^j(C)\,u^j,
\]
its generic cohomology character is
\begin{equation}\label{eq:alteredScalarUniformGraphs}
                    F(u)=\frac{d(u)}2(e+\tau).
\end{equation}
Every coefficient of \(d(u)\) is therefore even.

There is no such solution for which every \(H^j(X)\) is
graded flat as a left \(A\)-module. Thus a scalar-complex
repair would require nonflat cohomology. Neither the
explicit coupled Rouquier candidate nor the explicit
unit-cone direct-sum candidate from \fid{F59} can be
repaired by changing this scalar complex.
\end{finding}

\begin{proof}
\par\smallskip\noindent\textbf{Generic positive multiplicities.}\quad
Let \(Q\) be the ungraded fraction field of \(A\).
Ordinary generic localization is the semisimple graph
category with \(Q_g\otimes_Q Q_h\simeq Q_{gh}\) and
no maps between distinct graphs. The exact localization
framework is \cite[\S1.1, equations (1.1)--(1.2);
\S1.2, Theorem~1.6 and Remark~1.7]{EWLocalized}, with
the ordinary bimodule realization fixed in the earlier
\fid{F59} proof.

Forget internal grading at this stage. Write
\[
 F(u)=\sum_{g\in W}f_g(u)g,
 \qquad f_g(u)\in\mathbb N[u,u^{-1}],
\]
for the actual generic cohomological multiplicities.
K\"unneth over \(Q\) and the actual quadratic give
\begin{equation}\label{eq:alteredScalarGroupQuadratic}
                         F(u)^2=d(u)F(u).
\end{equation}
The ambient class and scalar Euler class imply
\begin{equation}\label{eq:alteredScalarGenericEuler}
                    F(-1)=2e+2\tau,\qquad d(-1)=4.
\end{equation}
Indeed, after forgetting internal shifts, \(b_s\) gives
\(e+s\) and \(b_s-\delta\) gives \(s-e\); their
two-factor expression is \(2(e+\tau)\).

Set \(a(u)=\sum_g f_g(u)\). Augmentation gives
\(a^2=da\). Since \(a\ne0\) and the Laurent ring is
an integral domain, \(a=d\). For each positive real
\(u\),
\[
                       \mu_u(g)=f_g(u)/d(u)
\]
is an idempotent probability measure on the finite group
\(W\). Its support \(S\) is independent of positive
\(u\), because each nonzero nonnegative Laurent polynomial
is positive there.

Here is the elementary finite-group argument. Positivity
in the convolution implies \(S^2=S\). A nonempty finite
subsemigroup of a group contains the identity and inverses,
so \(S\) is a subgroup. Choose \(g_0\in S\) maximizing
\(\mu(g)\). The identity
\[
 \mu(g_0)=\sum_{x\in S}\mu(x)\mu(x^{-1}g_0)
\]
expresses the maximum as a weighted average, with every
weight positive. Every term must attain that maximum.
Since \(x^{-1}g_0\) runs through \(S\), the measure is
uniform on \(S\).

Consequently \(f_g(u)=d(u)/|S|\) for \(g\in S\),
and it is zero otherwise. Equality for all positive real
\(u\) is a Laurent-polynomial identity. At \(u=-1\),
the coefficient of the identity in
\eqref{eq:alteredScalarGenericEuler} forces \(|S|=2\),
and its other element is \(\tau\). This proves
\eqref{eq:alteredScalarUniformGraphs}.

This conclusion retains only cohomological degree.
Generic localization inverts homogeneous polynomials;
in particular, a degree-two unit identifies internal
shifts differing by two. One cannot replace \(d(u)\)
here by the full internally graded character and infer
an immediate parity contradiction from \(\delta^2\).

\par\smallskip\noindent\textbf{Restoring the global grading under flatness.}\quad
Suppose every \(H^j(X)\) is graded left-flat.
These modules are finite, hence graded projective and
graded free. The last step follows by lifting a homogeneous
residue basis, splitting the graded free surjection,
and applying graded Nakayama to its complementary kernel,
as in the preceding proof.

Each cohomology module has generic support only \(e,\tau\).
Choose root coordinates \(x,y\) negated by \(\tau\)
and coordinates fixed by it. On a left-free cohomology
module, let \(U,V\) be right multiplication by \(x,y\).
Generically,
\[
                         yU=xV,\qquad U^2=x^2I.
\]
Left freeness extends these identities over \(A\).
Coprimality of \(x,y\) gives \(U=xZ,V=yZ\), with
\(Z\) homogeneous of degree zero and \(Z^2=I\).
The fixed coordinates act by their left scalars, also
by generic equality and torsion-freeness.

The bimodule idempotents \((I+Z)/2\) and \((I-Z)/2\)
split the module into globally defined \(e\)- and
\(\tau\)-graph pieces. Their left modules are graded
projective, hence graded free. They are therefore
right-free as well. There are actual nonnegative
Laurent multiplicity polynomials
\[
 a(v,u),b(v,u)\in\mathbb N[v,v^{-1},u,u^{-1}]
\]
such that the full global cohomology character is
\begin{equation}\label{eq:alteredScalarGlobalGraphs}
                       H_X(v,u)=a(v,u)e+b(v,u)\tau.
\end{equation}
Here \(u\) records cohomological degree, \(v\) records
actual global shifts, and \(b\ne0\) by
\eqref{eq:alteredScalarGenericEuler}.

All these cohomology modules are right-free. The canonical
cross-product map
\begin{equation}\label{eq:alteredScalarKunneth}
 \bigoplus_{i+j=r}H^i(X)\otimes_AH^j(X)
       \longrightarrow H^r(X\otimes_A^{\mathbf L}X)
\end{equation}
is therefore an isomorphism of graded bimodules. To verify
bijectivity, split the bounded complex as a complex of
right modules into its projective cohomology and
contractible pairs, and then tensor with the second
complex. These are only right-module splittings; they
prove bijectivity of the canonical cross-product map,
which is itself bimodule-linear. No bimodule formality
is inferred from a one-sided splitting.

Put
\[
 d(v,u)=\sum_{j,e}\dim_kH^j(C)_e\,v^e u^j.
\]
Scalar K\"unneth and the quadratic, with
\eqref{eq:alteredScalarKunneth}, give
\begin{equation}\label{eq:alteredScalarMultiplicityEquations}
                     a^2+b^2=da,\qquad 2ab=db.
\end{equation}
The graph types can be compared separately: a map
between distinct graphs would have its coefficient
killed by \(2x\) in the polynomial domain and is zero.
Since the integral Laurent ring is a domain and
\(b\ne0\), the second equation forces \(d=2a\).
At the cohomological Euler specialization \(u=-1\),
\[
                   v^2+2+v^{-2}=2a(v,-1),
\]
which is impossible because the two outer coefficients
on the left are odd. The full grading is used only
after the global free graph splitting has justified it.

\par\smallskip\noindent\textbf{The two explicit candidates.}\quad
The coupled Rouquier Euler lift has derived image
\[
 X_{\mathrm R}\simeq A^{\oplus2}
                      \oplus A_\tau(-2)\oplus A_\tau(2)
\]
in degree zero. If it satisfied the changed quadratic,
its square would be in degree zero. Since
\(H^j(C)\otimes_kX_{\mathrm R}\ne0\) whenever
\(H^j(C)\ne0\), the scalar complex would also have
cohomology only in degree zero. Its Euler class then
forces the original \(\delta^2\) coefficient.
But the identity-graph characters of the two sides are
respectively
\[
                  v^4+6+v^{-4},\qquad
                  2v^2+4+2v^{-2},
\]
which differ. Arbitrary scalar shifts do not repair it.

For the actual unit-dot cone, the earlier derived image
is \(C_s\simeq R_s(1)\oplus R(1)[1]\).
Consequently the candidate
\[
 X_{\mathrm{dot}}=(B_s\boxtimes B_s)
                       \oplus(C_s\boxtimes C_s)
\]
has ungraded generic cohomology character
\begin{equation}\label{eq:alteredScalarDotCharacter}
 F(z)=(1+z^2)e+(1+z)(s_1+s_2)+2\tau,
\end{equation}
where \(z\) records the shift \([1]\), hence
cohomological degree \(-1\). This cannot satisfy
\(F^2=dF\) for any scalar Laurent polynomial \(d\),
even without the prescribed Euler class.
The character evaluations on the four-element subgroup
generated by \(s_1,s_2\) give the three nonzero values
\[
                    z^2+2z+5,\qquad z^2-1,\qquad (z-1)^2.
\]
Each value \(f\ne0\) would satisfy \(f^2=df\),
forcing \(d=f\) in the Laurent polynomial domain.
The three values differ. Equivalently, its extra
\(s_1,s_2\) cohomology violates the uniform support
already proved. Neither argument retains any invalid
full internal grading after localization.
\end{proof}

\paragraph{What the finite Frobenius fiber still controls.}
For the central finite graded Frobenius quotient \(B=A/I\)
from Finding~\ref{f:boundedSoergelNoGo}, the equality
\(IM=MI\) on each term still gives the monoidal termwise
quotient. If the changed quadratic is an actual homotopy
quadratic, it descends for \(Z=X/IX\) to
\[
                     Z\otimes_B Z\simeq C\otimes_k Z.
\]
The Euler dimension of \(Z\) is \(4\dim_kB\ne0\),
so it has nonzero cohomology. Let \(r\) and \(a\)
be the upper cohomological degrees of \(Z\) and \(C\).
The local nonvanishing argument gives upper degree
\(2r\) on the left, while scalar K\"unneth gives
\(a+r\) on the right. Hence \(r=a\).
Exact graded Frobenius duality gives the corresponding
equality of the lower degrees. This use of termwise
quotient requires the stated homotopy isomorphism;
it is not a claim that ordinary quotient preserves
arbitrary derived isomorphisms.

The old amplitude argument now matches the extremes
of \(Z\) and \(C\), rather than forcing either into
degree zero. Its concentrated-projective-fiber and
minimal-complex conclusion cannot simply be repeated.
The remaining case is an actual bounded coupled lift
with nonflat cohomology and an altered scalar complex,
with actual relation maps. Neither a construction nor
an unconditional exclusion of that case is proved here.
Polynomial-differential coefficient bimodules and
unbounded completions are further changes of premise.

% Separate paragraph for Part III.
\paragraph{Altered scalar complexes in Grothendieck's language.}
The scalar \(C\) is a finite perfect complex of graded
vector spaces, within SGA~6, I, Definition~4.7,
pp.~121--122, and Example~4.8, p.~122 \cite{SGA6}.
Its Euler class is virtual, whereas the proof uses
actual nonnegative cohomological multiplicities.
The formal torus grading has the SGA~3, I, 4.7.2--3
framework and cote~90, scan~17; SGA~6, V, \S\S2--3,
and cote~133, scan~18, concern symmetric coefficient
operations \cite{SGA3,SGA6,Archive}.
No \(\lambda\)-ring structure on the noncommutative
Hecke word algebra is inferred.

The difference between localized graph multiplicities
and globally graded free summands is an affine module
and base-change distinction. EGA~I, Theorem~1.3.7,
p.~86, Corollary~1.3.8, p.~87, and Theorem~1.4.1,
pp.~90--92, supply the module/sheaf framework \cite{EGA}.
The K\"unneth and flatness implications needed here are
proved with their actual hypotheses. The elementary
finite-group uniform-measure argument is included as
background, not claimed as a new general theorem.
The new application is the cohomological-support and
nonflatness requirement for this generator and changed
scalar quadratic. The finite Frobenius algebra in
the proof carries no asserted arithmetic Frobenius,
Weil weight, positivity polarization or canonical basis.
